\documentclass{ltjbook}
\usepackage{docmute}
\usepackage{../../myStyle}

\begin{document}
\VerbatimFootnotes

\chapter{JASPでベイズ分析を体験する}
\label{ch28_bayesian_with_JASP}

% 章冒頭: GUIソフトウェアという選択肢の話をここに置く
% 移設元: ch27_bayes_judgement.old  L87--94「GUIソフトウェアの登場」
% 節見出しは立てず，既存の章冒頭リード文（ch28_bayesian_with_jasp.old L9付近）と統合して地の文にする

\section{JASPの導入とインターフェイス}
% 移設元: ch28_bayesian_with_jasp.old  L26--70「JASPの導入とインターフェイス」

\section{記述統計と可視化}
% 移設元: ch28_bayesian_with_jasp.old  L71--106「記述統計と可視化」

\section{JASPによる相関・回帰分析}
% 移設元: ch28_bayesian_with_jasp.old  L107--145「JASPによる相関・回帰分析」

\section{ベイズファクターを使った検定}
% 移設元: ch28_bayesian_with_jasp.old  L146--289「ベイズファクターを使った検定」

\end{document}
